\chapter{Quantum symmetries and geometric comparisons}
\label{ch:vol2-quantum-interfaces}

An algebra of boundary operators determines how those operators compose.
It does not determine how two representations combine.
The latter operation uses a coproduct, or a more general tensor
construction on the category of representations.
Nor does an exchange operator determine the coefficient lattice on
which multiplication acts. These distinctions persist in explicit
polynomial examples.

We begin with one algebra carrying two different Hopf structures.
We then compare ordinary bars, polynomial shuffle lattices and
matrix-coefficient reconstruction. The final examples concern
geometric and gravitational comparisons.
A tangent complex does not determine its nonlinear equations.
The spectrum of a connection does not identify its holonomy with
its coefficients.

Throughout, $k$ is a field of characteristic zero.
All tensors are algebraic until a different completion or localization
is specified.

\section{One operator algebra and two representation tensors}

Put $A=k[x,y,z]$, with augmentation
$\varepsilon(x)=\varepsilon(y)=\varepsilon(z)=0$.
For $\eta=0,1$, define algebra maps
$\Delta_\eta:A\to A\otimes_k A$ by
\begin{equation}\label{eq:vol2-two-hopf}
\begin{aligned}
 \Delta_\eta x&=x\otimes1+1\otimes x,\\
 \Delta_\eta y&=y\otimes1+1\otimes y,\\
 \Delta_\eta z&=z\otimes1+1\otimes z+\eta x\otimes y.
\end{aligned}
\end{equation}
A Hopf structure requires coassociativity, the two counit identities,
and an antipode $S$ with
$m(S\otimes1)\Delta=m(1\otimes S)\Delta=\mathbf1\varepsilon$.
Here $m$ is multiplication and $\mathbf1:k\to A$ is the unit.

\begin{proposition}\label{prop:vol2-two-hopf}
Both coproducts in \eqref{eq:vol2-two-hopf} define Hopf algebras.
Their antipodes are
\[
 S_\eta x=-x,\qquad S_\eta y=-y,\qquad
 S_\eta z=-z+\eta xy.
\]
They have the same augmented algebra and the same ordinary bar
coalgebra, but they are not isomorphic as Hopf algebras.
\end{proposition}
\begin{proof}
It suffices to check the generators.
The only nonprimitive generator is $z$.
Both iterated coproducts give its three single-factor terms and
\[
 \eta\bigl(x\otimes y\otimes1
       +x\otimes1\otimes y+1\otimes x\otimes y\bigr).
\]
This proves coassociativity.

The augmentation kills each extra tensor factor, giving the counits.
On $z$, either antipode composite is
$(-z+\eta xy)+z-\eta xy=0$.
The checks on $x,y$ are immediate, and the identities extend
multiplicatively.

The coproduct $\Delta_0$ is cocommutative, whereas
\[
 \Delta_1z-\Delta_1^{\mathrm{op}}z=x\otimes y-y\otimes x\ne0.
\]
Cocommutativity is preserved by a Hopf isomorphism, so none exists.
The ordinary augmented bar uses multiplication, augmentation and
unit, with word deconcatenation as its coalgebra structure.
Those data are identical for the two algebras.
\end{proof}

The difference is visible on one-dimensional modules.
Let $M$ have character $(x,y,z)=(1,0,0)$ and let $N$ have character
$(0,1,0)$.
On $M\otimes_k N$, the action induced by $\Delta_0$ has $z=0$.
The action induced by $\Delta_1$ has $z=1$.
Thus even the ordinary tensor of representations is extra information
relative to the augmented algebra and its bar.

Assigning weights $1,1,2$ to $x,y,z$ makes both coproducts graded.
Their common Hilbert series
\[
 \frac{1}{(1-t)^2(1-t^2)}
\]
does not distinguish them either.

\section{Ordinary bars and geometric coefficients}

Let $V$ be finite-dimensional and concentrated in degree zero.
Give its generators weight one.
The tensor algebra and the square-zero algebra are
\[
 T(V)=\bigoplus_{r\ge0}V^{\otimes r},
 \qquad
 A_{\mathrm{sq}}=k\oplus V,\qquad V^2=0.
\]
Both are ordinary augmented algebras over $k$.

We use a cohomological bar convention: a word of bar length $n$
has degree $-n$, and $sV=V[1]$ has degree $-1$.
For an augmentation ideal $\bar A$, the bar has terms
$(s\bar A)^{\otimes n}$ and differential given by adjacent products.
Its coproduct cuts a word at every position, including both empty cuts.

\begin{proposition}\label{prop:vol2-free-bar}
There is a quasi-isomorphism of coaugmented coalgebras
\[
 k\oplus sV\longrightarrow B(T(V)),
\]
where $sV$ is primitive.
Consequently the derived augmentation end algebra is
\[
 \operatorname{RHom}_{T(V)}(k,k)
 \simeq k\oplus V^*[-1]
\]
with square-zero reduced multiplication.
For the square-zero algebra,
\[
 B(A_{\mathrm{sq}})=T^c(sV),
\]
with zero differential. Its continuous weight dual is
$\widehat T(V^*[-1])=\prod_{n\ge0}(V^*[-1])^{\otimes n}$,
with concatenation product.
\end{proposition}
\begin{proof}
Choose a basis of $V$.
A bar element of $T(V)$ is a nonempty word in basis elements,
together with a set of cuts separating its nonempty algebra factors.
Fix an underlying word of length $r$.
Its possible cuts form subsets
$C\subset\{1,\ldots,r-1\}$, and its bar length is $|C|+1$.

Choose the bar sign for which
\[
 dC=\sum_{j=1}^{|C|}(-1)^{j-1}
       \bigl(C\setminus\{c_j\}\bigr),
 \qquad c_1<\cdots<c_{|C|}.
\]
For $r\ge2$, define $hC=\{1\}\cup C$ when $1\notin C$, and
$hC=0$ otherwise.

If $1\notin C$, the term deleting the new first cut is $C$.
The other terms cancel those of $hdC$.
If $1\in C$, only the term of $dC$ deleting that cut contributes
to $hdC$, again giving $C$.
Hence $dh+hd=1$ on this word.
Only words of length one and the empty word survive.
Their inclusion respects deconcatenation and proves the first assertion.

The ordinary two-sided free bar resolution computes the derived Hom.
Dualizing the displayed coalgebra quasi-isomorphism gives the
square-zero multiplication, since the surviving reduced classes are
primitive. Dualization is exact over $k$.
For $A_{\mathrm{sq}}$, every product of two augmentation-ideal
elements is zero, so the multiplication differential vanishes.
Each weight is finite-dimensional. Dualizing the algebraic weight sum
therefore gives the indicated product and the dual concatenation law.
\end{proof}

These formulas require an ordinary scalar augmentation.
An empty word, a local vacuum and a prepared state have different
definitions. None can replace that augmentation without a map.
Likewise, the coefficient dual
$(\bigoplus_n W_n)^*=\prod_n W_n^*$ changes a sum to a product.
It does not by itself identify a field dual, an internal line dual
or the boundary-junction construction of
Chapter~\ref{ch:vol2-marked-reconstruction}.

\subsection{Componentwise tensor algebras and factorization}

For a factorization coefficient $F$, restriction to two distinct
points identifies $F_{\{1,2\}}$ with
$F_{\{1\}}\otimes F_{\{2\}}$.
Suppose these three coefficients are $k$.
The componentwise sum $F\oplus F$ has coefficient $k^2$ at
both one and two points, whereas factorization would require $k^4$
at two points.

The same difficulty applies to the naive componentwise tensor algebra.
For the coefficient $F$ just described, it has $k[t]$ at every
nonempty configuration. The inherited separated map is
\[
 k[t]\longrightarrow k[t_1,t_2],
 \qquad t^n\longmapsto t_1^nt_2^n.
\]
It misses $t_1$ and is not a factorization equivalence.
Weight completion changes the map to
$k[[t]]\to k[[t_1,t_2]]$ with the same missing coefficient.

The free and square-zero calculation remains valid in its ordinary
algebraic category. Placing it in a factorization category requires
the actual free construction there, or a comparison functor that
preserves the required structure.
Pointwise tensor closure alone does not give closure under the
coproducts used to build a tensor algebra.

\begin{proposition}\label{prop:vol2-bar-functor}
Let $F:\mathcal C\to\mathcal D$ be a strong monoidal functor
between monoidal categories in which the augmented simplicial bars
and their geometric realizations exist.
Suppose $F$ preserves those realizations.
Then for every augmented associative algebra $A$ there is a natural
equivalence
\[
 F(B_{\mathcal C}A)\simeq B_{\mathcal D}(F A).
\]
It preserves the specified unit, augmentation and all simplicial maps.
\end{proposition}
\begin{proof}
In simplicial degree $n$, the strong monoidal structure identifies
$F(A^{\otimes n})$ with $(F A)^{\otimes n}$.
It carries adjacent-multiplication faces to faces, augmentation
faces to augmentation faces, and unit insertions to degeneracies.
The coherence axioms make these identifications a map of simplicial
objects. Apply preservation of geometric realization.
\end{proof}

An ordered collision atlas and a category of consecutive word cuts
have different diagrams.
Applying the proposition to them first requires a functor on those
diagrams and a comparison of their coefficient tensors.
It also requires preservation of any support condition and completion
used in the realization.
An ordinary bar--cobar equivalence on either carrier does not supply
this functor.
The Lie-to-commutative-coalgebra chiral duality has yet another operadic
source and target. It cannot replace the associative comparison.

\subsection{A finite Chevalley model is not a finite full bar}

For $A=U(\mathfrak{sl}_2)$, PBW gives the basis
$e^ah^bf^c$, excluding the constant when forming the augmentation
ideal $I$.
Thus $I^{\otimes n}$ is countably infinite-dimensional for every
$n\ge1$.
Its finite associated-graded pieces satisfy
\[
 \sum_{d\ge0}
 \dim\operatorname{gr}^{\mathrm{PBW}}_d(I^{\otimes n})\,t^d
 =\bigl((1-t)^{-3}-1\bigr)^n.
\]
The numbers $3^n$ occur only in the least possible PBW degree $d=n$.

The Chevalley complex instead has dimensions $1,3,3,1$.
With $d(x\wedge y)=[x,y]$,
\[
 d(e\wedge h)=-2e,\qquad
 d(e\wedge f)=h,\qquad
 d(h\wedge f)=-2f,
\]
while
\[
 d(e\wedge h\wedge f)
 =-2e\wedge f-h\wedge h-2f\wedge e=0.
\]
Its homology has dimensions $1,0,0,1$.
The bar--Chevalley comparison identifies homology, not these
different chain spaces.

\section{Zeros, poles and polynomial shuffle lattices}

Let $K(u)$ be a rational kernel.
For symmetric functions in $m$ and $n$ variables, sum over
the $\binom{m+n}{m}$ ordered two-block partitions.
Insert one factor $K$ for each cross-block pair.
This defines their shuffle product.
There is no factorial divisor.

An iterated product in three blocks has one summand for each ordered
three-block partition and the same cross factors in either bracketing.
This proves associativity wherever the stated coefficient space is
closed under the formula.

Consider first the meromorphic kernel
\[
 K_{\mathrm{mer}}(u)
 =\frac{(u+\epsilon_1)(u+\epsilon_2)(u+\epsilon_3)}{u^3},
 \qquad \epsilon_1+\epsilon_2+\epsilon_3=0.
\]
Put $\sigma_2=\sum_{i<j}\epsilon_i\epsilon_j$ and
$\sigma_3=\epsilon_1\epsilon_2\epsilon_3$.
Then
\begin{equation}\label{eq:vol2-shuffle-divisor}
 K_{\mathrm{mer}}(u)=1+\frac{\sigma_2}{u^2}
                         +\frac{\sigma_3}{u^3},
 \qquad
 K_{\mathrm{mer}}(u)-K_{\mathrm{mer}}(-u)
 =\frac{2\sigma_3}{u^3}.
\end{equation}
If the parameters are nonzero and pairwise distinct,
$u=-\epsilon_a$ is a zero, not a pole.
The reciprocal has residue
\begin{equation}\label{eq:vol2-shuffle-reciprocal}
 \operatorname{Res}_{u=-\epsilon_a}
    K_{\mathrm{mer}}(u)^{-1}\,du
 =\frac{(-\epsilon_a)^3}
        {\prod_{b\ne a}(\epsilon_b-\epsilon_a)}.
\end{equation}
Both statements follow by factoring the numerator at that point.

Degenerate parameter fibers must be calculated before cancelling
coincident zeros and poles.

For the convention using $K_{\mathrm{mer}}(z_1-z_2)$, put
$u=z_1-z_2$ and $w=z_2$.
The commutator of the degree-one inputs $z^r,z^s$ is
\[
 \frac{2\sigma_3}{u^3}
   \bigl((w+u)^rw^s-(w+u)^sw^r\bigr).
\]
For $(r,s)=(1,0)$ it is $2\sigma_3/u^2$.
For $(r,s)=(2,0)$ it is
$4\sigma_3w/u^2+2\sigma_3/u$.
The latter has residue $2\sigma_3$ at the true diagonal.

Neither expression has a pole at a generic nonzero shifted zero
$u=-\epsilon_a$.
To assign shifted residues to a geometric operation one must first
specify the meromorphic coefficient of that operation.

\subsection{The coefficient lattice retained by multiplication}

Now use the distinct polynomial Hall carriers
\begin{equation}\label{eq:vol2-hall-ring}
\begin{gathered}
 R=\mathbb Q[h_1,h_2],\qquad h_3=-h_1-h_2,\\
 a=-(h_1^2+h_1h_2+h_2^2),\qquad
 b=-h_1h_2(h_1+h_2),\\
 N(v)=v^3+av+b,\qquad
 K_0(v)=N(v)/v,\qquad K_1(v)=vN(v).
\end{gathered}
\end{equation}
For $j=0,1$, independently set
\[
 \mathcal A^{(j)}=\bigoplus_{n\ge0}A_n^{(j)},\qquad
 A_n^{(j)}=R[x_1,\ldots,x_n]^{\mathfrak S_n}.
\]
Every cochain degree and parity is zero.
Charge is the separate grading $n$. Each element has finite charge
support and polynomial coefficients.
No parameter or root is inverted.

The product is
\begin{equation}\label{eq:vol2-hall-shuffle}
 (F*_jG)(X)
 =\sum_{\substack{I\sqcup J=[d+e]\\ |I|=d}}
     F(X_I)G(X_J)
       \prod_{\substack{i\in I\\ k\in J}}K_j(x_k-x_i).
\end{equation}
The unit is $1\in A_0^{(j)}=R$.
The constant polynomial $e_0=1\in A_1^{(j)}$ has charge one and
is not that unit.

\begin{proposition}\label{prop:vol2-polynomial-shuffles}
Formula \eqref{eq:vol2-hall-shuffle} gives a polynomial,
unital associative algebra.
Both kernels have exchange ratio
\[
 r(v)=\frac{K_j(-v)}{K_j(v)}
     =\frac{v^3+av-b}{v^3+av+b}.
\]
\end{proposition}
\begin{proof}
For $K_1$ polynomiality is immediate.
For $K_0$, each possible denominator is a single root difference
in a summand.
Pair two summands by exchanging the labels of that pair between
the two blocks.
On their diagonal all other factors agree and the two residues
are opposite.
Every possible denominator therefore cancels in the polynomial
unique factorization domain over $R$.

Symmetry follows by permuting the summands.
The ordered-partition argument proves associativity in every charge.
An empty block proves the unit identities.
Finally $K_1=v^2K_0$, and this even factor cancels in the ratio.
\end{proof}

The ordinary zero-potential Hall product has precisely the
numerator Euler factors from moving arrows and denominator factors
from the flag tangent space \cite[Section~2.4, Theorem~2]{KSCoHA2011}.
Three moving arrow weights $v+h_i$ give $K_0$.
Adding two zero-weight loops multiplies this kernel by $v^2$.
This geometric interpretation does not identify these ordinary
polynomial algebras with a critical Hall algebra or a native field
theory.

In charge two put
\[
 c=(x+y)/2,\qquad v=y-x,\qquad t=v^2,\qquad B=R[c,t].
\]
The entire input tensor is $R[x,y]=B\oplus vB$.
For $p=A+vC$, its linearized product is
\begin{equation}\label{eq:vol2-hall-binary}
 \mu_0(p)=2\bigl((t+a)A+bC\bigr),\qquad
 \mu_1(p)=2t\bigl((t+a)A+bC\bigr).
\end{equation}
This follows by adding $K_j(v)p(x,y)$ and
$K_j(-v)p(y,x)$.
Hence the complete binary images are
$J=(t+a,b)B$ and $tJ$.

\begin{theorem}\label{thm:vol2-hall-lattices}
There is no charge-preserving $R$-linear algebra isomorphism
$\mathcal A^{(0)}\to\mathcal A^{(1)}$.
The conclusion still holds for maps semilinear under one automorphism
of $R$.
Nevertheless the maps
\begin{equation}\label{eq:vol2-discriminant-map}
 \Phi_n(F)=D_nF,\qquad
 D_n=\prod_{i<k}(x_k-x_i)^2,\qquad D_0=D_1=1,
\end{equation}
give a unital injective algebra map.
They identify the two subalgebras generated by charge one.
They also give an isomorphism after inverting $D_n$ in each
charge, and commute with simultaneous root translation.
\end{theorem}
\begin{proof}
An isomorphism would identify the cokernels of the two full
binary products.
But
\[
 B/J=R[c]/(b),\qquad
 B/tJ\twoheadrightarrow R[c],\quad F\longmapsto F(c,0).
\]
The first has parameter annihilator $(b)$ and the second has
annihilator zero.
These cannot be isomorphic as $R$-modules.
Under a coefficient automorphism, the nonzero annihilator becomes
$(\sigma(b))$, so the semilinear case has the same contradiction.

For the positive map, the discriminant identity
\[
 D_{d+e}(X)
 =D_d(X_I)D_e(X_J)
    \prod_{\substack{i\in I\\ k\in J}}(x_k-x_i)^2
\]
holds for every ordered partition.
Replace $K_1$ by $v^2K_0$ in each summand of
$\Phi_d(F)*_1\Phi_e(G)$.
The discriminants combine to $D_{d+e}$, independent of the
partition. This proves multiplicativity.

The polynomial rings are domains, so each $\Phi_n$ is injective.
It is the identity in charges zero and one.
Its image on the charge-one-generated subalgebra is exactly the
corresponding generated subalgebra. Hence the restricted map is
an isomorphism.
Its inverse divides by $D_n$ only on this specified image.

After inverting $D_n$, the same division is defined on the whole
charge-$n$ carrier.
Products remain in this localization: every internal and cross
root difference divides a power of the total discriminant.
Finally $\sum_i\partial_{x_i}D_n=0$, which proves translation
compatibility.
\end{proof}

This is a distinction between the full polynomial realization and
the smaller algebra generated by its currents.
Equality of exchange ratios decides neither.

The quotient
$Q_n=A_n^{(1)}/D_nA_n^{(1)}$ retains the discarded collision
coefficients.
It is an $\mathcal A^{(0)}$-bimodule through
\[
 a\cdot[z]=[\Phi(a)*_1z],\qquad
 [z]\cdot a=[z*_1\Phi(a)].
\]
These are well-defined because $\Phi$ is multiplicative.

The quotient is not an algebra quotient: the image of $\Phi$
contains the unit and is a proper subalgebra.
In charge two its cone is
$[B\xrightarrow{t}B]$, with cohomology $B/(t)$ in degree zero.
On the ordered cover this is the length-two normal quotient
$R[c,v]/(v^2)$.

\section{Cooperation depends on the tensor}

The product on a polynomial shuffle carrier does not automatically
admit a coproduct into the ordinary tensor product.
The obstruction uses only charge one and charge two.
Write $e_m=x^m\in A_1^{(j)}$.
Formula \eqref{eq:vol2-hall-binary} gives
\begin{equation}\label{eq:vol2-cooperation-relation}
 [e_0,e_3]_{*_j}-3[e_1,e_2]_{*_j}
       +a[e_0,e_1]_{*_j}
 =b(e_0*_je_0).
\end{equation}
Indeed, the three commutators on the left are
$2bt^j(x^2+xy+y^2)$, $2bt^jxy$ and $2bt^j$, while
$e_0*_je_0=2t^j(t+a)$.

\begin{theorem}\label{thm:vol2-shuffle-no-cooperation}
Neither $\mathcal A^{(j)}$ admits a charge-preserving unital
$R$-algebra map
\[
 \Delta:\mathcal A^{(j)}
       \longrightarrow\mathcal A^{(j)}\otimes_R\mathcal A^{(j)}
\]
with the ordinary tensor multiplication and with two-sided counit
equal to projection onto charge zero.
The same obstruction applies to the current-generated subalgebras,
including a derived algebra map with the ordinary Koszul tensor.
Coassociativity is not needed for this conclusion.
\end{theorem}
\begin{proof}
Charge preservation leaves only the splits $(1,0)$ and $(0,1)$
for each $e_m$.
The two counits force
$\Delta e_m=e_m\otimes1+1\otimes e_m$.
Commutators of even primitive elements are primitive, since their
mixed terms cancel.
Thus the reduced coproduct of the left side of
\eqref{eq:vol2-cooperation-relation} is zero.
Multiplicativity makes the reduced coproduct of its right side
\[
 2b\,e_0\otimes e_0\ne0
       \quad\text{in }A_1\otimes_R A_1=R[x,y].
\]
This is a contradiction over the original coefficient ring.

For the ambient algebras all charge components are $R$-free.
For a current-generated subalgebra, taking degree-zero cohomology
of a proposed derived map still gives the same identity in the
free $(1,1)$ component.
Hence the derived formulation has the same contradiction.
\end{proof}

The detector can vanish on a special parameter fiber.
The theorem asserts no uniform fibrewise obstruction when $b=0$.
At $(h_1,h_2,h_3)=(1,1,-2)$, however, it becomes
$-4e_0\otimes e_0$, so the obstruction is already visible over
$\mathbb Q$.
Neither this result nor its proof concerns every possible
cut-convolution coalgebra or every shifted coefficient tensor.

\subsection{A rational cooperation with a different interchange}

There is an explicit positive construction after localizing the
spectral coefficients and changing the middle interchange.
For ordered root lists $X,Y$, put
\[
 Q_{d,e}^{(j)}(X,Y)
   =\prod_{x\in X,\ y\in Y}K_j(y-x),\qquad
 \Delta_{d,e}^{\mathrm{rat}}F
   =\frac{F(X,Y)}{Q_{d,e}^{(j)}(X,Y)}.
\]
The recipient consists of block-symmetric polynomials localized
at all pair differences and their translates by $\pm h_i$.
Parameters themselves are not inverted.

Let $F$ have bidegree $(d,e)$ and $G$ have bidegree $(d',e')$.
Define their product by summing over
$X=X_1\sqcup X_2$, $Y=Y_1\sqcup Y_2$ of these sizes, with summand
\begin{equation}\label{eq:vol2-rational-interchange}
\begin{aligned}
 &F(X_1,Y_1)G(X_2,Y_2)
   \prod_{\substack{x\in X_1\\x'\in X_2}}K_j(x'-x)\\
 &\qquad\cdot
   \prod_{\substack{y\in Y_1\\y'\in Y_2}}K_j(y'-y)
   \prod_{\substack{y\in Y_1\\x'\in X_2}}r(y-x').
\end{aligned}
\end{equation}
The last product is the middle interchange.
Its factors would be absent from the ordinary tensor multiplication.

\begin{proposition}\label{prop:vol2-rational-cooperation}
The multiplication \eqref{eq:vol2-rational-interchange} is
associative and unital.
The sum of the maps $\Delta_{d,e}^{\mathrm{rat}}$ is multiplicative
into this recipient and is rationally coassociative.
It has both charge-zero counits.
\end{proposition}
\begin{proof}
For multiplicativity, consider a term with four ordered groups
$X_1,Y_1,X_2,Y_2$.
Divide its source shuffle product by the cross kernel between
$X_1\sqcup X_2$ and $Y_1\sqcup Y_2$.
The factors from $X_1$ to $Y_2$ cancel.
The remaining middle ratio is
\[
 \frac{K_j(x'-y)}{K_j(y-x')}=r(y-x'),
 \qquad y\in Y_1,\ x'\in X_2.
\]
The other factors are exactly those in
\eqref{eq:vol2-rational-interchange}, together with
$\Delta F$ and $\Delta G$.
Summing over all four groups proves the identity.

For three inputs, either bracketing includes the same internal
kernel factors and the same factors from an earlier right block
to a later left block.
This proves associativity.
For three output lists, either iterated cooperation divides by
$Q_{12}Q_{13}Q_{23}$ in the common localization.
This proves coassociativity.
The empty lists give the unit and both counits.
Every input has finite charge support, so these sums are finite.
\end{proof}

The four-current relation is now compatible with
\[
 (1-r(v))v(v^2+a)=b(1+r(v)),
\]
which follows directly from
$r(v)=(v^3+av-b)/(v^3+av+b)$.
The mixed contribution excluded by the ordinary tensor is present.
The construction still leaves the polynomial carrier:
\[
 \Delta_{1,1}^{\mathrm{rat}}(1\in A_2)
  =\frac{v}{N(v)}\quad(j=0),\qquad
 \Delta_{1,1}^{\mathrm{rat}}(1\in A_2)
  =\frac{1}{vN(v)}\quad(j=1).
\]
It also removes the closed loci at its inverted denominators.
An equivalence on this open coefficient domain is not yet a
supported comparison across those loci.

\subsection{Shifted collisions retain their excess}

Keep the same coefficient ring $R$.
On three position variables consider
\[
 f_{12}=z_1-z_2+h_1,\qquad
 f_{23}=z_2-z_3+h_1.
\]
Their intersection is $R[z_3]$, with
$z_2=z_3-h_1$ and $z_1=z_3-2h_1$.
The composed endpoint shift is $-2h_1$.
Thus the finite list of shifts $0,\pm h_1,\pm h_2,\pm h_3$
is not closed under composition of these shifted collisions.

For an ordered list of equations, use the Koszul complex with
degree-minus-one generators $\xi_{ij}$ and
$d\xi_{ij}=f_{ij}$.
Adding the induced endpoint equation gives
\[
 f_{13}=z_1-z_3+2h_1=f_{12}+f_{23}.
\]
The change of exterior generator
$\eta=\xi_{13}-\xi_{12}-\xi_{23}$ yields
\begin{equation}\label{eq:vol2-shifted-excess}
 K(f_{12},f_{23},f_{13})
 \simeq R[z_3]\otimes\Lambda(\eta),
 \qquad |\eta|=-1,\qquad d\eta=0.
\end{equation}
Indeed, $f_{12},f_{23},z_3$ are polynomial coordinates over $R$,
so their first two Koszul pairs contract. $\eta$ remains.
Discarding the redundant equation would discard its nonzero
negative-degree class.

The true endpoint diagonal instead restricts to $-2h_1$.
Its derived restriction is
\[
 [R[z_3]\xrightarrow{-2h_1}R[z_3]]
\]
in degrees $-1,0$.
Before specialization its only cohomology is $R[z_3]/(h_1)$
in degree zero.
Derived specialization to $h_1=0$ makes the differential zero and
retains both degrees.

There is a separate obstruction to viewing the scalar quotient as
the unchanged supported differential module.
For $P=R[z_1,z_2]$ and $f=z_1-z_2-\lambda$, with
$\lambda\in R$, a relative differential-module structure on
$P/(f)$ would satisfy $[\partial_{z_1},f]=1$.
But $f$ acts by zero on that module, so its commutator with every
operator is zero.
The quotient is nonzero, giving a contradiction.

The same argument applies to a differential-module complex whose
underlying $P$-complex is quasi-isomorphic to $P/(f)$.
Its degree-zero cohomology would have that impossible action.
A supported differential-module realization must retain its normal
derivatives and the maps attaching them.
The scalar Koszul calculation does not supply those maps.

\section{Matrix coefficients reconstruct a specified tensor category}

The preceding examples distinguish a product, its cooperation and
its coefficient domain.
Tensor reconstruction begins with all of them encoded in a category.
Let $\mathcal C$ be an essentially small $k$-linear monoidal category,
and let
\[
 \omega:\mathcal C\longrightarrow\operatorname{Vect}_{\mathrm{fd}}(k)
\]
be a $k$-linear strong monoidal functor.
Write $V_M=\omega(M)$ and use its monoidal identifications
$V_M\otimes V_N\simeq V_{M\otimes N}$ and $V_{\mathbf1}\simeq k$.
Choose a small skeleton when forming the direct sum below.

The matrix-coefficient coend is the vector space
\begin{equation}\label{eq:vol2-coend}
 H_\omega=
 \left(\bigoplus_{M\in\mathcal C}V_M^*\otimes V_M\right)
 \big/
 \left\langle
 c_M(\lambda\omega(f),v)-c_N(\lambda,\omega(f)v)
 \right\rangle.
\end{equation}
Here $f:M\to N$, $v\in V_M$, $\lambda\in V_N^*$, and
$c_M(\lambda,v)$ denotes the class of the indicated tensor.
The relations remember every morphism of $\mathcal C$.

\begin{proposition}\label{prop:vol2-coend-operations}
The formulas
\begin{equation}\label{eq:vol2-coend-coalgebra}
\begin{aligned}
 \Delta c_M(\lambda,v)
   &=\sum_i c_M(\lambda,e_i)\otimes c_M(e^i,v),\\
 \varepsilon c_M(\lambda,v)&=\lambda(v)
\end{aligned}
\end{equation}
are independent of a basis $(e_i)$ and its dual $(e^i)$.
They define a counital coassociative coalgebra.
Its universal coaction is
\[
 \delta_M(v)=\sum_i e_i\otimes c_M(e^i,v).
\]
The monoidal structure supplies the multiplication
\begin{equation}\label{eq:vol2-coend-product}
 c_M(\lambda,v)c_N(\mu,w)
   =c_{M\otimes N}(\lambda\otimes\mu,v\otimes w),
\end{equation}
with unit $c_{\mathbf1}(1,1)$.
Together these operations make $H_\omega$ a bialgebra.
If $\mathcal C$ is rigid, it has an invertible antipode.
\end{proposition}
\begin{proof}
The tensor $\sum_i e_i\otimes e^i$ is the identity of $V_M$,
so its insertion in \eqref{eq:vol2-coend-coalgebra} is
basis-independent.
For $f:M\to N$, inserting this identity before or after $\omega(f)$
gives the same linear map.
The coend relation therefore makes $\Delta$ well-defined.
The same relation respects evaluation, so it also defines
$\varepsilon$.

Inserting the identity twice gives
\[
 \sum_{i,j}
 c_M(\lambda,e_i)\otimes
 c_M(e^i,e_j)\otimes
 c_M(e^j,v)
\]
for either iterated coproduct.
Evaluating either outer pair gives $c_M(\lambda,v)$, proving
both counits.
These same calculations prove the coaction identities.
The coend relation says exactly that each $\omega(f)$ is
colinear.

Naturality of the monoidal constraints makes
\eqref{eq:vol2-coend-product} compatible with the relations.
Their associativity and unit coherences give the algebra laws.
Using the basis $e_i\otimes f_j$ of $V_M\otimes V_N$ shows
that $\Delta$ preserves multiplication. Evaluation shows the
same for the counit.

For a left dual $M^\vee$, identify $V_{M^\vee}$ with $V_M^*$.
Let $\widetilde v\in V_M^{**}$ be evaluation at $v$.
Then
\[
 S\,c_M(\lambda,v)=c_{M^\vee}(\widetilde v,\lambda).
\]
The evaluation morphism $M^\vee\otimes M\to\mathbf1$ and the
coevaluation morphism $\mathbf1\to M\otimes M^\vee$ are
morphisms of $\mathcal C$.
Applying their coend relations gives the two antipode identities:
the two sums contract to $\lambda(v)c_{\mathbf1}(1,1)$.
Right duals give the inverse antipode by the corresponding
right-dual formula.
\end{proof}

Coevaluation gives the coproduct, whereas the tensor product of
objects gives multiplication.
Composition has a third role.
The full linear dual $H_\omega^*$ is the algebra of linear natural
endomorphisms of $\omega$.
An endomorphism $T$ gives the functional
$c_M(\lambda,v)\mapsto\lambda(T_Mv)$, and
\[
 (\phi_T*\phi_U)c_M(\lambda,v)
   =\sum_i\lambda(T_Me_i)e^i(U_Mv)
   =\lambda(T_MU_Mv).
\]
Thus convolution is composition.

Tensor-compatible natural automorphisms must additionally preserve
the monoidal constraints and the unit.
They are not all elements of this dual algebra.

\begin{theorem}[Tensor reconstruction]\label{thm:vol2-tensor-reconstruction}
Suppose $\mathcal C$ is $k$-linear abelian, with
finite-dimensional morphism spaces and finite-length objects.
Suppose it is rigid, its tensor product is exact in both variables,
and $\omega$ is exact and faithful.
Then the comparison given by $\delta_M$ is a monoidal equivalence
\[
 \mathcal C\simeq
       \operatorname{Comod}_{\mathrm{fd}}(H_\omega).
\]
\end{theorem}
\begin{proof}
The coalgebra reconstruction theorem for an exact faithful functor
to finite-dimensional vector spaces identifies $\mathcal C$ with
the finite-dimensional right comodules of the coend
\eqref{eq:vol2-coend}
\cite[Theorem~1.10.1]{EGNO2015}.
Its comparison is exactly the coaction
\eqref{eq:vol2-coend-coalgebra}.
Equation \eqref{eq:vol2-coend-product} identifies the coaction
on $V_{M\otimes N}$ with the tensor product of the two coactions.
The given monoidal constraints therefore make the equivalence
strong monoidal, including its unit.
Proposition~\ref{prop:vol2-coend-operations} supplies the antipode,
as in the Hopf reconstruction theorem
\cite[Theorem~5.4.1]{EGNO2015}.
\end{proof}

For a braided $\mathcal C$, write
$\beta_{M,N}:V_M\otimes V_N\to V_N\otimes V_M$ for the induced
linear exchange.
Its coefficient functional is
\begin{equation}\label{eq:vol2-coend-braiding}
 \mathfrak r\bigl(c_M(\lambda,v),c_N(\mu,w)\bigr)
   =(\mu\otimes\lambda)\beta_{M,N}(v\otimes w).
\end{equation}
Naturality makes this well-defined on both coends.
The inverse braiding gives its convolution inverse.

Taking matrix coefficients of the two hexagon diagrams gives the
two coquasitriangular identities. Taking coefficients of the
colinearity square gives compatibility with multiplication.
Thus \eqref{eq:vol2-coend-braiding} is the coquasitriangular form.

Under the reconstruction equivalence it recovers the braiding by
\[
 v\otimes w\longmapsto
    \sum w_{(0)}\otimes v_{(0)}
      \mathfrak r(v_{(1)},w_{(1)}),
 \qquad
 \delta(v)=\sum v_{(0)}\otimes v_{(1)}.
\]

When $H_\omega$ is infinite-dimensional, the coefficient functional
lives in $(H_\omega\otimes H_\omega)^*$.
Regard $H_\omega^*$ with the inverse-limit topology from its finite
coefficient coalgebras, and define
$H_\omega^*\widehat\otimes H_\omega^*
=(H_\omega\otimes H_\omega)^*$ with that topology.
The functional is then the corresponding completed $R$-operator.
This does not place it in the ordinary algebraic tensor
$H_\omega^*\otimes H_\omega^*$.
A meromorphic version additionally needs specified expansion maps,
domains and continuity for its spectral parameter.

\subsection{An exchange relation is not a presentation}

For a selected object, the naturality of its exchange gives the
matrix relation usually written
$RT_1T_2=T_2T_1R$.
These relations define an algebra with a comparison map to the coend.
They form a presentation precisely when that map is an isomorphism.
The category's morphisms can impose further relations.

For example, compare finite-dimensional vector spaces with
finite-dimensional $\mathbb Z$-graded vector spaces, using the
ordinary flip and the forgetful fiber functor in both cases.
For vector spaces, every map $k\to V$ is allowed.
The coend relations therefore give
\[
 c_V(\lambda,v)=\lambda(v)c_k(1,1),\qquad H_\omega=k.
\]

For graded vector spaces, decompose into homogeneous lines.
There are no morphisms between distinct degrees.
There is one coefficient $u_n$ in each degree, and the formulas give
\[
 \Delta u_n=u_n\otimes u_n,\qquad
 u_nu_m=u_{n+m},\qquad u_0=1,\qquad S(u_n)=u_{-n}.
\]
Hence $H_\omega=k[t,t^{-1}]$ with $u_n=t^n$.

The exchange on a chosen one-dimensional fiber is the scalar $1$
in both examples.
The one-by-one RTT relation imposes no further condition on its
coefficient $t$.
In the first category the morphisms additionally impose $t=1$.
The same scalar exchange therefore leads to different reconstructed
Hopf algebras.

The geometric action algebra of a collection of kernels records
the operators they generate on a state space.
The coend instead records natural symmetries respecting every
chosen kernel morphism.
Identifying these two constructions requires a comparison of their
actions and tensor structures.
A list of root multiplicities, or a common exchange ratio, provides
neither that map nor its inverse.

\section{Geometric comparison preserves equations, not only counts}

Two positively graded Lie algebras can have the same graded PBW
character and different brackets.
Give $x,y$ weight one and $z$ weight two.
The abelian Lie algebra on these generators and the Heisenberg
Lie algebra with $[x,y]=z$ have the same PBW Hilbert series
$((1-t)^2(1-t^2))^{-1}$.
One enveloping algebra is commutative and the other is not.
The equality of their characters therefore supplies no algebra
isomorphism.

Similarly, the zero super vector space and a space with one even
and one odd basis vector both have superdimension zero.
An index cannot recover the forgotten parity dimensions.

An equivalence of underlying tangent complexes is also weaker
than a comparison of deformation equations.
Let $L$ have basis $u$ in degree one and $v$ in degree two, with
zero differential.
Consider the two dg Lie structures
\[
 L_{\mathrm{ab}}:\ [\, ,\,]=0,\qquad
 L_{\mathrm{q}}:\ [u,u]=2v,\quad [u,v]=[v,v]=0.
\]
The graded bracket on the odd element $u$ is symmetric.
Every iterated bracket is zero because $v$ is central, so both
structures satisfy Jacobi.
Their tangent complexes are identical.

\begin{proposition}\label{prop:vol2-nonlinear-comparison}
The first-order deformation $\epsilon u$ over
$k[\epsilon]/(\epsilon^2)$ extends to
$k[\epsilon]/(\epsilon^3)$ for $L_{\mathrm{ab}}$ and does not
extend for $L_{\mathrm{q}}$.
Thus the identity of tangent complexes cannot extend to an
$L_\infty$ equivalence between these deformation problems.
\end{proposition}
\begin{proof}
The Maurer--Cartan equation is
$d\alpha+\tfrac12[\alpha,\alpha]=0$.
It vanishes identically in the abelian case.
In the second algebra every possible lift has the form
$\alpha=(\epsilon+c\epsilon^2)u$.
Its curvature modulo $\epsilon^3$ is $\epsilon^2v\ne0$.
There is no degree-zero element giving a gauge correction.

Equivalently, an $L_\infty$ map with identity first component
would have to preserve the binary bracket up to the differential
of its second component.
Both differentials are zero, so it would identify the zero
bracket with $[u,u]=2v$, which is impossible.
\end{proof}

A geometric field comparison must consequently carry the nonlinear
operations as well as its tangent map.
For a cotangent BV theory it must also carry the pairing, action
and boundary term.
The classical master equation of a supplied coupled algebra proves
the consistency of that algebra.
It does not construct the action of an independently specified
gravitational or matter field theory on it.
At the quantum level, the chosen regularization and the local
quantum master equation remain additional comparisons, as in
Chapter~\ref{ch:vol2-native-bulk-boundary}.

\subsection{A projector must respect the actual transition maps}

A finite decomposition model has a simple idempotent:
on a sum indexed by pairs $(d_1,d_2)$, retain precisely the
summands with $d_1=d_2$.
Writing its inclusion and retraction as $\iota,\pi$ gives
\[
 \pi\iota=1,\qquad \iota\pi=e_{\mathrm{al}}.
\]
It is an equivalence on the full object exactly when the
off-aligned summands vanish.
These equations determine the invertibility of this canonical
comparison. They do not rule out every possible noninvertible
distributive law.

To descend the projector, equality of indices must be compatible
with each transition correspondence.
Preserving equality alone is insufficient.
The map $\{0,1\}\to\{*\}$ preserves equal pairs but sends the
off-diagonal pair $(0,1)$ to the diagonal pair $(*,*)$.
Extension along this map does not commute with the two projectors.

In a six-functor construction, assume the actual squares with the
diagonal index locus are Cartesian.
Assume also that their base-change maps are isomorphisms.
Under those hypotheses, base change identifies the pull--push
idempotents across every transition. The triple-diagonal identity
then descends as well.
Finite combinatorics alone does not verify those geometric squares.

\section{Gravitational spectra, holonomy and boundary data}

Three-dimensional Cartan gravity begins with independent fields
$e,\omega$.
For negative cosmological constant, the two connections are
$A^\pm=\omega\pm\ell^{-1}e$.
On the invertible-coframe locus, their flatness equations give
torsion zero and constant curvature.
This is a field construction with an action comparison.
Specifying two boundary central charges or a list of current spins
does not supply these fields or their action on the boundary.

\subsection{The principal embedding and its spin spectrum}

Let $V_m=\operatorname{Sym}^m(k^2)$ be the irreducible
$\mathfrak{sl}_2$ module of highest weight $m$.
The principal embedding in $\mathfrak{sl}_N$ acts on its defining
representation as $V_{N-1}$.
Its adjoint therefore decomposes as
\begin{equation}\label{eq:vol2-principal-spectrum}
 \mathfrak{sl}_N
 \cong\bigoplus_{s=2}^{N}V_{2s-2},\qquad
 \sum_{s=2}^{N}(2s-1)=N^2-1.
\end{equation}
Indeed, the tensor product
$V_{N-1}\otimes V_{N-1}^*$ contains once each highest weight
$0,2,\ldots,2N-2$.

This follows by subtracting successively the characters of those
highest weights from its weight character. Finite-dimensional
$\mathfrak{sl}_2$ modules are semisimple.
Removing the scalar endomorphisms removes $V_0$, proving
\eqref{eq:vol2-principal-spectrum}.

For a specified principal Chern--Simons theory linearized about
its gravitational $\mathfrak{sl}_2$ background, these summands
are the frame-like fields of spins $2,\ldots,N$.
With principal Drinfeld--Sokolov boundary conditions, the reduced
classical currents have the same spins
\cite{CampoleoniHigherSpin,HenneauxRey}.
This matches the representations in that construction.
Arbitrary supplied $W_N$ boundary data need a comparison from this
boundary reduction.
It must preserve the Poisson operations and the bulk coupling.
The dimension identity alone does not construct it.

\subsection{Connection coefficients and holonomy are different functions}

Work over $\mathbb C$ in this subsection, and let $L>0$.
Consider a constant traceless companion matrix $A(q)$ of size $N$:
\[
 A_{i,i+1}=1\quad(1\le i<N),\qquad
 A_{N,j}=q_{N+1-j}\quad(1\le j<N),\qquad A_{N,N}=0,
\]
with all remaining entries zero.
These are the constant coefficients of the corresponding
companion-gauge oper.
Expansion of the characteristic determinant gives
\begin{equation}\label{eq:vol2-companion-polynomial}
 \det(\lambda I-A(q))
   =\lambda^N-\sum_{s=2}^{N}q_s\lambda^{N-s}.
\end{equation}
One can also recover the same polynomial from the eigenvector
recursion $v_{i+1}=\lambda v_i$ and the final row.

For a circle of length $L$ and the parallel-transport convention
$\partial_\theta\psi=A(q)\psi$, the holonomy is
$H(q)=\exp(LA(q))$.
It is an exponential of the matrix, not its characteristic
polynomial.

\begin{proposition}\label{prop:vol2-holonomy-logarithm}
Holonomy conjugacy does not in general determine the coefficients
$q_s$.
Suppose $A(q)$ is regular semisimple.
Choose a domain containing its eigenvalues on which
$\lambda\mapsto e^{L\lambda}$ is injective.
The corresponding logarithm branch then recovers its coefficients
from holonomy conjugacy.
\end{proposition}
\begin{proof}
For $N=2$,
\[
 A(q)=\begin{pmatrix}0&1\\q&0\end{pmatrix},\qquad A(q)^2=qI,
\]
so the exponential series gives
\begin{equation}\label{eq:vol2-sl2-holonomy}
 H(q)=\cosh(L\sqrt q)\,I
       +\frac{\sinh(L\sqrt q)}{\sqrt q}\,A(q),
 \qquad
 \operatorname{Tr}H(q)=2\cosh(L\sqrt q).
\end{equation}
Both coefficients are entire functions of $q$. The second has
value $L$ at $q=0$.
For $L=2\pi$ and $q=-n^2$, every positive integer $n$ gives
$H(q)=I$.
Distinct connection coefficients therefore have the same holonomy.
This is already a counterexample among constant companion
connections.

Under the additional hypothesis, the distinct eigenvalues of
$H(q)$ give the unordered eigenvalues
$\lambda_i=L^{-1}\log\lambda_i(H)$ of $A(q)$.
They determine its characteristic polynomial.
Equation \eqref{eq:vol2-companion-polynomial} then recovers every
$q_s$.
\end{proof}

For nonconstant boundary fields the holonomy is a path-ordered
exponential. Its traces are global functions of those fields.
The local Drinfeld--Sokolov current brackets are obtained from
Hamiltonian reduction. They are not identified with the algebra
of these holonomy traces by a matrix characteristic polynomial.
Boundary conditions, a chosen logarithm or spectral family, and
the actual Hamiltonian comparison are therefore part of any
charge-identification theorem.
The constant example does not decide which of its connections
satisfies a separately imposed physical regularity condition.

\subsection{Algebraic removal of a central curvature}

Let $\mathfrak g$ be a dg Lie algebra and let
$\kappa\in\mathfrak g^2$ be closed and central.
Adjoin a central generator $u$ of degree one with $du=\kappa$.
The differential remains square-zero and a derivation:
the only new derivation equation is $[\kappa,x]=0$.
The class of $\kappa$ becomes exact in the larger algebra.
If $\kappa$ is the curvature term of the corresponding curved
Maurer--Cartan equation, twisting by $-u$ changes that term to
\[
 \kappa+d(-u)+\tfrac12[-u,-u]=0.
\]
This is an explicit algebraic construction.

It can change the deformation problem: for
$\mathfrak g=k\kappa$ with zero differential, the original
degree-two cohomology is nonzero, whereas the extension
$[ku\xrightarrow{d}k\kappa]$ is acyclic.

A geometric anomaly cancellation requires more.
Suppose an independently constructed invertible field theory has
the opposite anomaly.
A comparison must identify both classes in the same complex and
normalization.
Under these hypotheses, tensoring the theories cancels their sum.
Adjoining $u$ does not construct that invertible theory, its
descent data or the comparison of anomaly complexes.
The algebraic curvature, a determinant-line curvature and a
local BV obstruction retain their separate meanings until
such maps are supplied.

\section{The remaining geometric comparisons}

The constructions above preserve definite operations.
Ordinary bars retain augmented multiplication.
The discriminant map retains the shuffle product and translation
on its stated image.
Rational cooperation retains a changed middle interchange on an
open coefficient domain.
Tensor reconstruction retains a specified category, fiber
functor and all its morphisms.
Cartan gravity and companion-matrix holonomy begin with actual
connection fields.

To compare any of these with native chiral quantum symmetry, the
map must be defined before its equivalence is asserted.
For the shuffle models it must retain the integral binary images,
the shifted collisions and the excess class
\eqref{eq:vol2-shifted-excess}, with a specified tensor and counit.
For boundary reconstruction it must retain the derived local
mapping objects and relative sewing of
Problem~\ref{prob:vol2-marked-physical-reconstruction}.
For a gravitational realization it must carry the nonlinear
field equations, pairing, action and boundary reduction,
followed by the chosen quantum and state comparisons.
These requirements leave room for restricted coefficient domains,
larger mapping categories and interfaces between different
realizations. Each possibility has its own comparison map and
its own test of invertibility.
